\documentclass[10pt]{article}
\usepackage[margin=.75in]{geometry}
\usepackage{enumitem}
\usepackage{booktabs}
\usepackage{xcolor}
\usepackage{listings}
\usepackage{hyperref}

\setlength{\parskip}{.4em}
\setlength{\parindent}{0pt}

\lstset{
  language=Python,
  basicstyle=\ttfamily\small,
  frame=single,
  breaklines=true,
  showstringspaces=false
}

\title{\textbf{Assignment 1}\\
       Building a Reliable Mini Data Pipeline in Jupyter}
\author{}
\date{Fri Aug 28 2026}

\begin{document}
\maketitle

\section*{Scenario}
You are a junior data engineer supporting an e-commerce analytics pipeline.
A CSV file containing customer events arrives from several systems. Before
the data can be used by analysts, it must be validated, cleaned, transformed,
and checked for performance problems.

The file \texttt{events.csv} intentionally contains realistic
data-quality problems such as duplicate records, missing identifiers,
malformed timestamps, inconsistent text categories, and invalid numeric values.

\section*{Learning Goals}
By completing this assignment, you will practice:
\begin{itemize}
    \item writing reusable Python functions;
    \item debugging errors and interpreting tracebacks in Jupyter;
    \item building basic data-quality checks;
    \item timing and profiling Python code;
    \item organizing a small ETL-style workflow.
\end{itemize}

\section*{Dataset}
The dataset contains the following columns:

\begin{center}
\begin{tabular}{ll}
\toprule
Column & Meaning \\
\midrule
\texttt{event\_id} & Unique identifier for an event \\
\texttt{event\_timestamp} & Time at which the event occurred \\
\texttt{customer\_id} & Customer identifier \\
\texttt{product\_id} & Product identifier \\
\texttt{event\_type} & view, add\_to\_cart, purchase, or refund \\
\texttt{quantity} & Number of units \\
\texttt{unit\_price} & Price per unit \\
\texttt{country} & Customer country code \\
\texttt{source} & web, mobile, or partner\_api \\
\texttt{status} & Processing status \\
\bottomrule
\end{tabular}
\end{center}

\section*{Tasks}

\subsection*{1. Build reusable cleaning functions (3 pts)}
Load the CSV into a Jupyter notebook using pandas. Create and use at least
\textbf{three functions}. Your functions should collectively:

\begin{enumerate}[label=\alph*)]
    \item standardize \texttt{event\_type} and \texttt{status} by trimming
    whitespace and converting text to a consistent case;
    \item convert \texttt{event\_timestamp}, \texttt{quantity}, and
    \texttt{unit\_price} to appropriate data types;
    \item remove exact duplicates and identify invalid or missing values.
\end{enumerate}

For invalid numeric or timestamp values, convert them to missing values
rather than allowing the entire pipeline to fail.

\textbf{Rubric:} 1 point for function design, 1 point for correct cleaning,
and 1 point for clear function calls/docstrings or comments.

\subsection*{2. Debug a broken transformation (3 pts)}
Start with the following intentionally broken function:

\begin{lstlisting}
def calculate_revenue(df):
    df["revenue"] = df["quantity"] * df["unit_price"]
    purchases = df[df["event_type"] == "Purchase"]
    return purchases.groupby("country")["reveneu"].sum()
\end{lstlisting}

Run it on your cleaned data and use the traceback to debug it.

\begin{enumerate}[label=\alph*)]
    \item Explain at least \textbf{two problems} in the original function.
    \item Correct the function.
    \item Show the resulting purchase revenue by country.
\end{enumerate}

Use at least one Jupyter/IPython debugging feature such as
\texttt{\%xmode Context}, \texttt{\%xmode Verbose}, or \texttt{\%debug}.
Include a short Markdown note describing how the traceback helped you.

\textbf{Rubric:} 1 point for identifying problems, 1 point for a working
correction, and 1 point for demonstrating/debugging explanation.

\subsection*{3. Profile two implementations (2 pts)}
Create two implementations that calculate revenue by country:

\begin{itemize}
    \item \textbf{Version A:} a Python loop over records or rows;
    \item \textbf{Version B:} a pandas/vectorized solution using filtering,
    column arithmetic, and \texttt{groupby}.
\end{itemize}

Use \texttt{\%timeit} to compare them. Then use \texttt{\%prun} on at least
one implementation.

Answer in 2--4 sentences:
\begin{itemize}
    \item Which version is faster?
    \item What did the profiler show?
    \item Why does profiling matter in a data-engineering pipeline?
\end{itemize}

\textbf{Rubric:} 1 point for correct timing/profiling and 1 point for
interpretation.

\subsection*{4. Produce a pipeline quality summary (2 pts)}
Create a final function called:

\begin{lstlisting}
def pipeline_summary(raw_df, clean_df):
    ...
\end{lstlisting}

It should return or display at least:

\begin{itemize}
    \item number of raw rows;
    \item number of cleaned rows;
    \item number of duplicate rows removed;
    \item number of records with unusable timestamps;
    \item number of records with missing customer or product IDs;
    \item total valid purchase revenue.
\end{itemize}

Add a final Markdown paragraph explaining whether you would allow this
dataset to continue to a production analytics table and why.

\textbf{Rubric:} 1 point for the summary metrics and 1 point for the
data-engineering judgment.

\section*{Submission}
Submit one Jupyter notebook (\texttt{.ipynb}) containing:
\begin{itemize}
    \item executable code;
    \item outputs from the major steps;
    \item timing/profiling output;
    \item short Markdown explanations;
    \item a final cleaned-data and pipeline-quality summary.
\end{itemize}

Your notebook should run from top to bottom without manual fixes.

%\section*{Dataset Generation}
%The provided dataset was generated synthetically with a fixed random seed
%so that it is reproducible. The generator first creates valid e-commerce
%events, then deliberately injects data-quality problems and duplicate rows.
%
%To regenerate the dataset, run:
%\begin{lstlisting}
%python events.py
%\end{lstlisting}
%
%In a Jupyter notebook, you may also run:
%\begin{lstlisting}
%%run events.py
%\end{lstlisting}

\section*{Point Summary}
\begin{center}
\begin{tabular}{lr}
\toprule
Task & Points \\
\midrule
Reusable cleaning functions & 3 \\
Debugging a broken transformation & 3 \\
Timing and profiling & 2 \\
Pipeline quality summary & 2 \\
\midrule
\textbf{Total} & \textbf{10} \\
\bottomrule
\end{tabular}
\end{center}

\end{document}
